\section{Introduction}
\label{sec:intro}

Cloud back ends are assembled from fine-grained remote procedure calls across dependency graphs~\cite{dean2013tailatscale,gan2019deathstarbench}. A request completes only when its slowest required branch completes, so a modest queueing delay at one shared downstream service can dominate the end-to-end tail~\cite{dean2013tailatscale}. Production traces show that these graphs are deep and skewed, with a few heavily shared services carrying most of the traffic~\cite{luo2021characterizing}, and the study of their faults has become a field of its own~\cite{zhou2021fault}. Cloud hosts are heterogeneous and shared, so the latency of one service depends on the work that runs beside it~\cite{reiss2012heterogeneity,delimitrou2013paragon}. Outages of large cloud services are frequent enough to be studied systematically~\cite{gunawi2016cloud}. Measurements of single-server tail latency attribute much of the tail to interference and scheduling rather than to the application itself~\cite{li2014tales,fried2020caladan}, and warehouse-scale profiling shows that these effects are pervasive~\cite{kanev2015profiling}.

Thread-pool concurrency adds a second source of coupling that is easy to overlook. A synchronous upstream handler holds its worker thread for the full duration of each downstream call. An upstream pool can therefore saturate while the host processors are still mostly idle, because the scarce resource is the thread, not the CPU. The pool then accumulates queueing delay that belongs to a downstream tier. Metastability studies describe how this kind of overload feedback keeps a system degraded after the original trigger has passed~\cite{bronson2021metastable,huang2022metastable}. Admission-side controllers address overload at a single service, such as priority-based shedding~\cite{zhou2018dagor} or credit-based admission for microsecond-scale RPCs~\cite{cho2020breakwater}. Neither accounts for the thread-holding that a deadline budget must cross hop by hop. Sidecar proxies add a further per-hop cost, which has been measured for production-style service meshes~\cite{zhu2023sidecar}.

The problem this paper addresses is how these two couplings combine as call graphs grow. A chain of synchronous hops and a wide in-cast of shared leaves saturate at different loads and for different reasons, and the operator needs to know which resource binds first, how far the load can rise before tails explode, and whether the latency cost of observability and proxying is material at those loads. These questions require controlled experiments under open-loop load, because closed-loop generators understate saturation by throttling their own arrivals when servers slow down~\cite{schroeder2006open}.

We make four contributions. The first is a coupled holding-time model for bounded call graphs (Section~\ref{sec:model}). It predicts the root pool utilization from the arrival rate and the longest holding path, and it yields estimators for tail amplification and for a saturation horizon. The second is a process-per-vertex gRPC testbed with three concurrency runtimes, open-loop Poisson load with coordinated-omission correction, and a quiet-host protocol that rejects runs disturbed by background load (Section~\ref{sec:testbed}). The third is a campaign of 540 completed runs in five phases that tests the model on chains, scatter-gather trees, diamond in-casts, and two DeathStarBench-derived sub-graphs, and that measures the serialization and sidecar tax, the tracing overhead, and the behavior of three overload controllers under three injected faults (Sections~\ref{sec:campaign} and~\ref{sec:results}). The fourth is DAC-Shed, a deadline-budget and latency-gradient controller that is compared with a static circuit breaker and with naive retries, together with an offline replay of a budget-gated single retry.

The main results are as follows. P99.9 latency at the highest stable load was 2.8 to 8.1 times the critical-path median sum for 17 of the 18 Phase~1 shapes, and it was 20.5 times for a width-eight in-cast whose shared leaves approach saturation. Along chains, the saturation horizon fell from about 312 requests per second at depth two to about 77 at depth eight, while the root pool reached utilization 0.75 to 0.79 and the host processes used about ten per cent of the eight cores. At depth eight, synchronous chains timed out on 11.7\% of requests at 115 requests per second, while asynchronous and hybrid chains of the same depth showed no timeouts up to 191 requests per second. At 77 and 115 requests per second the three runtimes used similar CPU, and above that the synchronous CPU stops rising because its requests fail. The model over-predicted measured stable-load P99.9 by about 1.6 times at the median. The chain error traces to the root's fixed shape parameter, and two added diagnostics reduced the chain error but did not remove the fan-out error. Under a one-second hang or a four-fold slowdown of the third vertex at high load, DAC-Shed kept the P99.9 of successful requests at or below 256\,ms, where the static breaker and naive retries reached several seconds. Naive retries delivered the most goodput under transient errors, and an offline replay of a budget-gated single retry recovered 69.2\% of first-attempt errors. Tracing with sampling rates of 0.25 or higher raised median P99 by 4.3 to 10.6\,ms at high load (pooled median +5.50\,ms, $p = 0.03125$); because the rate-zero baseline ran under higher background CPU (0.93 vs.\ 0.11--0.41 cores), this shift provides a conservative lower bound on high-load tracing overhead.

Section~\ref{sec:threats} discusses the scope and validity boundaries of our single-host process-per-vertex testbed, bytes-passthrough sidecar emulation, and trace-replay counterfactual. A multi-phase sensitivity analysis of the quiet-host gate (Section~\ref{sec:gate}) confirms that relaxing the acceptance threshold from 1.0 to 1.2 cores leaves all 18 Phase~1 saturation brackets and all Phase~3--5 conclusions unchanged.

The remainder of the paper is organized as follows. Section~\ref{sec:related} places the work against the literature. Section~\ref{sec:model} states the model and the estimators, and Section~\ref{sec:testbed} describes the testbed, the runtimes, the controllers, and the measurement protocol. Section~\ref{sec:campaign} describes the campaign design and the quiet-host discards. Section~\ref{sec:results} reports the results for the five research questions and for the model diagnostics. Section~\ref{sec:gate} tests the sensitivity of the results to the quiet-host gate. Section~\ref{sec:threats} discusses the threats to validity, Section~\ref{sec:guidelines} gives deployment guidance, and Section~\ref{sec:conclusion} concludes.

The research questions that organize the evaluation are these. RQ1 asks how far high-quantile latency amplifies relative to the critical-path median, and whether the amplification depends on depth. RQ2 asks where saturation occurs, in the thread pool or on the processors, and how the three concurrency runtimes differ. RQ3 asks how large the serialization and sidecar costs are relative to CPU and latency. RQ4 asks how much distributed tracing adds to tail latency near saturation. RQ5 asks whether deadline-budget shedding protects successful requests under injected faults, and how it compares with a static breaker and with naive retries.
